\chapter{Fiches de tests fonctionnels}

Cette annexe présente les fiches de tests fonctionnels prévues pour valider les principaux parcours de l'application \emph{5 Secondes Chrono}. 
Les scénarios sont construits à partir des personae définis dans le projet afin de vérifier que l'application répond aux besoins concrets des utilisateurs cibles.

Chaque fiche de test permet de vérifier un parcours ou une fonctionnalité importante, aussi bien du côté de l'application mobile apprenant que de l'interface administrateur.

\section{Méthodologie de test}

Les tests sont organisés autour de scénarios utilisateurs. Chaque scénario est associé à un persona, un objectif, des préconditions, des étapes de test et des critères de validation.

Les statuts utilisés sont les suivants :

\begin{itemize}
    \item \textbf{Non testé} : le test n'a pas encore été réalisé ;
    \item \textbf{Validé} : le résultat obtenu correspond au résultat attendu ;
    \item \textbf{Partiellement validé} : le parcours fonctionne, mais une anomalie ou une amélioration est identifiée ;
    \item \textbf{Non validé} : le résultat obtenu ne correspond pas au comportement attendu.
\end{itemize}

\section{Modèle de fiche de test}

\begin{table}[H]
\centering
\renewcommand{\arraystretch}{1.3}
\small
\begin{tabular}{|p{4cm}|p{10cm}|}
\hline
\rowcolor{membrepurple!80}
\multicolumn{2}{|c|}{\textbf{Fiche de test -- Modèle}} \\
\hline
\textbf{Identifiant du test} & À compléter \\ \hline
\textbf{Titre du test} & À compléter \\ \hline
\textbf{Persona concerné} & À compléter \\ \hline
\textbf{Interface testée} & Mobile apprenant / Interface administrateur \\ \hline
\textbf{Objectif du test} & À compléter \\ \hline
\textbf{Préconditions} & À compléter \\ \hline
\textbf{Données de test} & À compléter \\ \hline
\textbf{Résultat attendu} & À compléter \\ \hline
\textbf{Résultat obtenu} & À compléter après exécution du test \\ \hline
\textbf{Statut} & Non testé / Validé / Partiellement validé / Non validé \\ \hline
\textbf{Commentaires / anomalies} & À compléter \\ \hline
\end{tabular}
\caption{Modèle de fiche de test fonctionnel}
\end{table}

\section{Fiches de tests -- Application mobile apprenant}

\subsection{Test M01 -- Réaliser le quiz quotidien}

\begin{table}[H]
\centering
\renewcommand{\arraystretch}{1.3}
\small
\begin{tabular}{|p{4cm}|p{10cm}|}
\hline
\rowcolor{BleuPetrole!25}
\multicolumn{2}{|c|}{\textbf{Fiche de test M01 -- Quiz quotidien}} \\
\hline
\textbf{Identifiant du test} & M01 \\ \hline
\textbf{Titre du test} & Réaliser le quiz quotidien \\ \hline
\textbf{Persona concerné} & Nicolas, négociateur en immobilier commercial \\ \hline
\textbf{Interface testée} & Application mobile apprenant \\ \hline
\textbf{Objectif du test} & Vérifier qu'un apprenant peut lancer un quiz quotidien, répondre à une question chronométrée et consulter son résultat. \\ \hline
\textbf{Préconditions} & 
\begin{itemize}
    \item L'utilisateur possède un compte actif.
    \item L'utilisateur est connecté à l'application.
    \item Au moins une question est disponible dans la base de données.
\end{itemize}
\\ \hline
\textbf{Données de test} & 
Compte apprenant : à compléter. \newline
Question test : à compléter. \\ \hline
\textbf{Étapes de test} &
\begin{enumerate}
    \item Ouvrir l'application mobile.
    \item Se connecter avec un compte apprenant valide.
    \item Accéder à l'écran d'accueil.
    \item Lancer le quiz du jour.
    \item Répondre à la question proposée avant la fin du timer.
    \item Valider la réponse.
    \item Consulter le résultat affiché.
\end{enumerate}
\\ \hline
\textbf{Résultat attendu} & 
Le quiz se lance correctement. Le timer est visible. La réponse est prise en compte. Le score et le feedback sont affichés à l'utilisateur. \\ \hline
\textbf{Résultat obtenu} & À compléter après exécution du test. \\ \hline
\textbf{Statut} & Non testé / Validé / Partiellement validé / Non validé \\ \hline
\textbf{Commentaires / anomalies} & À compléter. \\ \hline
\end{tabular}
\caption{Fiche de test M01 -- Quiz quotidien}
\end{table}

\subsection{Test M02 -- Consulter une recommandation après une difficulté}

\begin{table}[H]
\centering
\renewcommand{\arraystretch}{1.3}
\small
\begin{tabular}{|p{4cm}|p{10cm}|}
\hline
\rowcolor{BleuPetrole!25}
\multicolumn{2}{|c|}{\textbf{Fiche de test M02 -- Recommandation personnalisée}} \\
\hline
\textbf{Identifiant du test} & M02 \\ \hline
\textbf{Titre du test} & Consulter une recommandation après une difficulté \\ \hline
\textbf{Persona concerné} & Sarah, jeune professionnelle en prise de poste \\ \hline
\textbf{Interface testée} & Application mobile apprenant \\ \hline
\textbf{Objectif du test} & Vérifier que l'application propose une recommandation simple lorsqu'un utilisateur rencontre une difficulté sur une thématique. \\ \hline
\textbf{Préconditions} &
\begin{itemize}
    \item L'utilisateur est connecté.
    \item L'utilisateur a réalisé au moins un quiz.
    \item Une erreur ou une difficulté est identifiable sur une thématique.
\end{itemize}
\\ \hline
\textbf{Données de test} & 
Compte apprenant : à compléter. \newline
Thématique testée : à compléter. \\ \hline
\textbf{Étapes de test} &
\begin{enumerate}
    \item Ouvrir l'application.
    \item Accéder à l'écran de progression ou de résultats.
    \item Identifier une thématique avec un score faible.
    \item Consulter la recommandation proposée.
    \item Ouvrir le contenu conseillé, par exemple une vidéo ou un quiz ciblé.
\end{enumerate}
\\ \hline
\textbf{Résultat attendu} & 
L'application affiche une recommandation compréhensible et liée à la difficulté détectée. L'utilisateur peut accéder au contenu conseillé. \\ \hline
\textbf{Résultat obtenu} & À compléter après exécution du test. \\ \hline
\textbf{Statut} & Non testé / Validé / Partiellement validé / Non validé \\ \hline
\textbf{Commentaires / anomalies} & À compléter. \\ \hline
\end{tabular}
\caption{Fiche de test M02 -- Recommandation personnalisée}
\end{table}

\subsection{Test M03 -- Reprendre rapidement l'application après une période d'inactivité}

\begin{table}[H]
\centering
\renewcommand{\arraystretch}{1.3}
\small
\begin{tabular}{|p{4cm}|p{10cm}|}
\hline
\rowcolor{BleuPetrole!25}
\multicolumn{2}{|c|}{\textbf{Fiche de test M03 -- Reprise rapide}} \\
\hline
\textbf{Identifiant du test} & M03 \\ \hline
\textbf{Titre du test} & Reprendre rapidement l'application après une période d'inactivité \\ \hline
\textbf{Persona concerné} & Karim, indépendant avec un agenda chargé \\ \hline
\textbf{Interface testée} & Application mobile apprenant \\ \hline
\textbf{Objectif du test} & Vérifier qu'un utilisateur irrégulier peut reprendre l'application facilement sans devoir reconfigurer son parcours. \\ \hline
\textbf{Préconditions} &
\begin{itemize}
    \item L'utilisateur possède un compte existant.
    \item L'utilisateur a déjà réalisé au moins un quiz.
    \item L'utilisateur n'a pas utilisé l'application depuis plusieurs jours.
\end{itemize}
\\ \hline
\textbf{Données de test} & 
Compte apprenant : à compléter. \newline
Dernière connexion simulée : à compléter. \\ \hline
\textbf{Étapes de test} &
\begin{enumerate}
    \item Ouvrir l'application après plusieurs jours d'inactivité.
    \item Observer l'écran d'accueil.
    \item Vérifier la présence d'un accès rapide au quiz.
    \item Lancer une session courte.
    \item Consulter le résultat et le feedback.
\end{enumerate}
\\ \hline
\textbf{Résultat attendu} & 
L'utilisateur retrouve rapidement son parcours. L'accès au quiz est visible. La session peut être réalisée sans configuration complexe. \\ \hline
\textbf{Résultat obtenu} & À compléter après exécution du test. \\ \hline
\textbf{Statut} & Non testé / Validé / Partiellement validé / Non validé \\ \hline
\textbf{Commentaires / anomalies} & À compléter. \\ \hline
\end{tabular}
\caption{Fiche de test M03 -- Reprise rapide}
\end{table}

\section{Fiches de tests -- Interface administrateur}

\subsection{Test A01 -- Créer une nouvelle question}

\begin{table}[H]
\centering
\renewcommand{\arraystretch}{1.3}
\small
\begin{tabular}{|p{4cm}|p{10cm}|}
\hline
\rowcolor{orange!20}
\multicolumn{2}{|c|}{\textbf{Fiche de test A01 -- Création d'une question}} \\
\hline
\textbf{Identifiant du test} & A01 \\ \hline
\textbf{Titre du test} & Créer une nouvelle question depuis l'espace administrateur \\ \hline
\textbf{Persona concerné} & Claire, administratrice / formatrice côté Le Carré Pro \\ \hline
\textbf{Interface testée} & Interface administrateur web \\ \hline
\textbf{Objectif du test} & Vérifier qu'un administrateur peut créer une question de quiz complète et l'enregistrer dans le système. \\ \hline
\textbf{Préconditions} &
\begin{itemize}
    \item L'administrateur possède un compte valide.
    \item L'administrateur est connecté à l'espace d'administration.
    \item La section de gestion des questions est accessible.
\end{itemize}
\\ \hline
\textbf{Données de test} & 
Question : à compléter. \newline
Réponses proposées : à compléter. \newline
Bonne réponse : à compléter. \newline
Explication : à compléter. \\ \hline
\textbf{Étapes de test} &
\begin{enumerate}
    \item Se connecter à l'espace administrateur.
    \item Accéder à la section de gestion des questions.
    \item Cliquer sur l'action de création d'une nouvelle question.
    \item Renseigner le texte de la question.
    \item Renseigner les réponses possibles.
    \item Définir la bonne réponse.
    \item Ajouter une explication.
    \item Enregistrer la question.
\end{enumerate}
\\ \hline
\textbf{Résultat attendu} & 
La question est enregistrée sans erreur. Elle apparaît dans la liste des questions et peut être utilisée dans un quiz. \\ \hline
\textbf{Résultat obtenu} & À compléter après exécution du test. \\ \hline
\textbf{Statut} & Non testé / Validé / Partiellement validé / Non validé \\ \hline
\textbf{Commentaires / anomalies} & À compléter. \\ \hline
\end{tabular}
\caption{Fiche de test A01 -- Création d'une question}
\end{table}


% ============================================================
% Fiche de test A01 -- Création d'une question
% ============================================================

\renewcommand{\arraystretch}{1.3}
\setlength{\LTleft}{\fill}
\setlength{\LTright}{\fill}

\begin{longtable}{
    |>{\raggedright\arraybackslash}p{0.28\textwidth}
    |>{\raggedright\arraybackslash}p{0.64\textwidth}|
}

\caption{Fiche de test A01 -- Création d'une question, complétée le 13/07/2026}
\label{tab:test-a01-creation-question}
\\

\hline
\rowcolor{orange!50}
\multicolumn{2}{|c|}{
    \textbf{Fiche de test A01 -- Création d'une question}
}
\\
\hline

\endfirsthead

\multicolumn{2}{c}{
    \small\textit{Fiche de test A01 -- Création d'une question (suite)}
}
\\
\hline

\rowcolor{orange!50}
\multicolumn{2}{|c|}{
    \textbf{Fiche de test A01 -- Création d'une question}
}
\\
\hline

\endhead

\hline
\multicolumn{2}{|r|}{
    \small\textit{Suite de la fiche page suivante}
}
\\
\hline

\endfoot

\hline
\endlastfoot

% ------------------------------------------------------------
% Informations générales
% ------------------------------------------------------------

\textbf{Identifiant du test}
&
A01
\\
\hline

\textbf{Titre du test}
&
Créer une nouvelle question depuis l'espace administrateur.
\\
\hline

\textbf{Persona concerné}
&
Claire, administratrice et formatrice pour Le Carré Pro.
\\
\hline

\textbf{Interface testée}
&
Interface web d'administration.
\\
\hline

\textbf{Date d'exécution}
&
13 juillet 2026.
\\
\hline

\textbf{Objectif du test}
&
Vérifier qu'un administrateur peut créer une question de quiz complète,
renseigner ses différentes propriétés, l'enregistrer dans le système et la
retrouver dans la liste des questions.
\\
\hline

% ------------------------------------------------------------
% Préconditions
% ------------------------------------------------------------

\textbf{Préconditions}
&
\begin{itemize}[leftmargin=*,nosep]
    \item L'administrateur possède un compte valide.
    \item Le compte utilisé dispose des droits d'administration.
    \item L'administrateur peut accéder à l'espace d'administration.
    \item La rubrique de gestion des questions est disponible.
    \item L'application et la base de données sont accessibles.
\end{itemize}
\\
\hline

% ------------------------------------------------------------
% Données de test
% ------------------------------------------------------------

\rowcolor{orange!15}
\multicolumn{2}{|c|}{
    \textbf{Données utilisées pour le test}
}
\\
\hline

\textbf{Type de bail attendu}
&
Bail commercial.
\\
\hline

\textbf{Thématique attendue}
&
État des lieux.
\\
\hline

\textbf{Question}
&
Quel est l'objectif principal d'un état des lieux d'entrée ?
\\
\hline

\textbf{Réponses proposées}
&
\begin{itemize}[leftmargin=*,nosep]
    \item \textbf{A --} Calculer la valeur de vente du local.
    \item \textbf{B --} Décrire précisément l'état du local au début de la location.
    \item \textbf{C --} Déterminer le chiffre d'affaires du locataire.
    \item \textbf{D --} Fixer automatiquement le montant du loyer.
\end{itemize}
\\
\hline

\textbf{Bonne réponse}
&
B -- Décrire précisément l'état du local au début de la location.
\\
\hline

\textbf{Explication pédagogique}
&
L'état des lieux d'entrée permet de décrire précisément l'état du local et
de ses équipements au début de la location. Il pourra ensuite être comparé
avec l'état des lieux de sortie afin d'identifier les éventuelles
dégradations intervenues pendant la durée du bail.
\\
\hline

\textbf{Statut de publication attendu}
&
Question active ou publiée, lorsque cette option est disponible dans le
formulaire.
\\
\hline

% ------------------------------------------------------------
% Étapes du test
% ------------------------------------------------------------

\rowcolor{orange!15}
\multicolumn{2}{|c|}{
    \textbf{Déroulement du test}
}
\\
\hline

\textbf{Étape 1}
&
Se connecter à l'espace administrateur avec un compte possédant le rôle
administrateur.
\\
\hline

\textbf{Étape 2}
&
Vérifier que la connexion aboutit bien au tableau de bord de
l'administration.
\\
\hline

\textbf{Étape 3}
&
Accéder à la rubrique \textit{Questions} ou
\textit{Gestion des questions}.
\\
\hline

\textbf{Étape 4}
&
Noter le nombre de questions déjà présentes dans la liste avant la création.
\\
\hline

\textbf{Étape 5}
&
Cliquer sur le bouton permettant d'ajouter ou de créer une nouvelle
question.
\\
\hline

\textbf{Étape 6}
&
Rechercher et sélectionner le type de bail
\textit{Bail commercial}.
\\
\hline

\textbf{Étape 7}
&
Rechercher et sélectionner la thématique
\textit{État des lieux}.
\\
\hline

\textbf{Étape 8}
&
Saisir la question suivante :
\textit{« Quel est l'objectif principal d'un état des lieux d'entrée ? »}
\\
\hline

\textbf{Étape 9}
&
Renseigner les quatre réponses proposées dans les champs correspondants.
\\
\hline

\textbf{Étape 10}
&
Définir la réponse B comme réponse correcte.
\\
\hline

\textbf{Étape 11}
&
Ajouter l'explication pédagogique prévue dans les données de test.
\\
\hline

\textbf{Étape 12}
&
Examiner le champ permettant d'associer une vidéo à la question et vérifier
la cohérence des vidéos proposées dans la liste déroulante.
\\
\hline

\textbf{Étape 13}
&
Activer ou publier la question lorsque cette option est proposée.
\\
\hline

\textbf{Étape 14}
&
Cliquer sur le bouton \textit{Enregistrer}, \textit{Créer} ou
\textit{Valider}.
\\
\hline

\textbf{Étape 15}
&
Vérifier qu'un message de confirmation est affiché et qu'aucun message
d'erreur bloquant n'apparaît.
\\
\hline

\textbf{Étape 16}
&
Rechercher la nouvelle question dans la liste des questions.
\\
\hline

\textbf{Étape 17}
&
Ouvrir la fiche de la question et vérifier que le texte, les réponses,
la bonne réponse, l'explication et les informations de classement ont été
correctement enregistrés.
\\
\hline

\textbf{Étape 18}
&
Actualiser la page, puis vérifier que la question est toujours présente.
\\
\hline

\textbf{Étape 19}
&
Utiliser la fonction de prévisualisation ou lancer un quiz de test afin de
vérifier l'affichage de la question et de ses réponses.
\\
\hline

\textbf{Étape 20}
&
Vérifier que la réponse B est reconnue comme correcte et que l'explication
associée est affichée après la réponse.
\\
\hline

\textbf{Étape 21}
&
Vérifier le fonctionnement du lien vers la vidéo éventuellement associée.
\\
\hline

\textbf{Étape 22}
&
Réaliser une ou plusieurs captures d'écran de la question créée et de son
affichage dans l'interface.
\\
\hline

\textbf{Étape 23}
&
À la fin du test, supprimer ou désactiver la question afin de ne pas
conserver une donnée de recette parmi les contenus définitifs.
\\
\hline

% ------------------------------------------------------------
% Résultats
% ------------------------------------------------------------

\rowcolor{orange!15}
\multicolumn{2}{|c|}{
    \textbf{Résultats du test}
}
\\
\hline

\textbf{Résultat attendu}
&
\begin{itemize}[leftmargin=*,nosep]
    \item Le formulaire de création est accessible.
    \item Les types de baux et les thématiques sont présentés avec des
    intitulés compréhensibles.
    \item Le type \textit{Bail commercial} peut être sélectionné.
    \item La thématique \textit{État des lieux} peut être sélectionnée.
    \item Les réponses et l'explication peuvent être enregistrées.
    \item La bonne réponse peut être clairement identifiée.
    \item Les vidéos disponibles sont correctement proposées dans la liste.
    \item La question est enregistrée sans erreur.
    \item La question apparaît dans la liste après validation.
    \item Les informations sont conservées après actualisation de la page.
    \item La question peut être utilisée et prévisualisée dans un quiz.
\end{itemize}
\\
\hline

\textbf{Résultat obtenu}
&
\begin{itemize}[leftmargin=*,nosep]
    \item Le formulaire de création d'une question est accessible depuis
    l'espace administrateur.

    \item La question a pu être renseignée et enregistrée dans
    l'application.

    \item Le type de bail \textit{Bail commercial} n'était pas disponible
    dans la liste proposée. La valeur \textit{Partie au contrat} a donc été
    utilisée pour poursuivre le test.

    \item La thématique \textit{État des lieux} n'était pas disponible.
    La valeur \textit{Professionnelle} a donc été sélectionnée à la place.

    \item Plusieurs intitulés affichés dans les menus correspondent
    directement aux noms techniques utilisés dans la base de données.
    Certains contiennent notamment des caractères de soulignement
    (\texttt{\_}).

    \item Le menu déroulant d'association d'une vidéo indiquait qu'aucune
    vidéo n'était disponible.

    \item Après l'enregistrement de la question, un lien vers une vidéo
    associée était pourtant affiché et ce lien fonctionnait correctement.

    \item La question créée était toujours présente après sa validation.

    \item La session administrateur est restée active après la fermeture
    de l'application et le redémarrage de l'ordinateur.
\end{itemize}
\\
\hline

\textbf{Statut}
&
\textbf{Partiellement validé}
\\
\hline

% ------------------------------------------------------------
% Anomalies constatées
% ------------------------------------------------------------

\rowcolor{orange!15}
\multicolumn{2}{|c|}{
    \textbf{Commentaires et anomalies constatées}
}
\\
\hline

\textbf{Anomalie A01-01}\newline
\textbf{Persistance de la session}
&
La session administrateur reste active après la fermeture du navigateur et
le redémarrage de l'ordinateur.

Aucune option visible ne permet à l'utilisateur d'indiquer s'il souhaite ou
non rester connecté.

Il conviendrait de vérifier la durée de validité du jeton
d'authentification, son stockage et son expiration. Une case
\textit{« Se souvenir de moi »} pourrait être ajoutée afin de laisser ce
choix à l'utilisateur. Le bouton de déconnexion doit également invalider
correctement la session.
\\
\hline

\textbf{Anomalie A01-02}\newline
\textbf{Intitulés techniques}
&
Plusieurs valeurs affichées dans les menus semblent provenir directement
des noms enregistrés dans la base de données.
Les caractères de soulignement et les intitulés techniques doivent être
remplacés par des libellés lisibles et compréhensibles pour
l'administrateur.

Par exemple, une valeur telle que
\texttt{bail\_commercial} devrait être affichée sous la forme
\textit{Bail commercial}.
\\
\hline

\textbf{Anomalie A01-03}\newline
\textbf{Catégories manquantes}
&
Les valeurs attendues \textit{Bail commercial} et
\textit{État des lieux} ne sont pas disponibles dans les listes proposées.

Cette absence oblige l'administrateur à sélectionner des catégories qui ne
correspondent pas réellement au contenu de la question. Cela peut fausser
le classement, la recherche et l'utilisation ultérieure de la question
dans les quiz.

Il est nécessaire de vérifier le contenu des tables de référence et les
valeurs chargées dans les menus déroulants.
\\
\hline

\textbf{Anomalie A01-04}\newline
\textbf{Association des vidéos}
&
Le menu déroulant indique qu'aucune vidéo n'est disponible, alors qu'un lien
vers une vidéo fonctionnelle apparaît après l'enregistrement de la question.
Le comportement du champ est donc incohérent avec les données réellement
enregistrées.

Il convient de vérifier :
\begin{itemize}[leftmargin=*,nosep]
    \item la requête utilisée pour charger la liste des vidéos ;
    \item la valeur sélectionnée par défaut ;
    \item l'identifiant de la vidéo enregistré avec la question ;
    \item le rafraîchissement du menu après la création ou la modification
    d'une vidéo ;
    \item l'utilité et le comportement attendu de ce champ.
\end{itemize}
\\
\hline

\textbf{Conclusion du test}
&
La fonction principale de création d'une question est opérationnelle :
le formulaire est accessible et la question peut être enregistrée.

Le test reste toutefois partiellement validé en raison de plusieurs
incohérences concernant les catégories proposées, les libellés affichés,
la gestion de l'association vidéo et la persistance de la session
administrateur.
\\
\hline

\end{longtable}



\subsection{Test A02 -- Associer une vidéo à une question ou une thématique}

\begin{table}[H]
\centering
\renewcommand{\arraystretch}{1.3}
\small
\begin{tabular}{|p{4cm}|p{10cm}|}
\hline
\rowcolor{orange!20}
\multicolumn{2}{|c|}{\textbf{Fiche de test A02 -- Association d'une vidéo}} \\
\hline
\textbf{Identifiant du test} & A02 \\ \hline
\textbf{Titre du test} & Associer une vidéo pédagogique à une question ou une thématique \\ \hline
\textbf{Persona concerné} & Claire, administratrice / formatrice côté Le Carré Pro \\ \hline
\textbf{Interface testée} & Interface administrateur web \\ \hline
\textbf{Objectif du test} & Vérifier qu'un administrateur peut rattacher une vidéo pédagogique à un contenu existant. \\ \hline
\textbf{Préconditions} &
\begin{itemize}
    \item L'administrateur est connecté.
    \item Une question ou une thématique existe déjà.
    \item Une vidéo est disponible ou peut être ajoutée.
\end{itemize}
\\ \hline
\textbf{Données de test} & 
Titre de la vidéo : à compléter. \newline
URL ou fichier vidéo : à compléter. \newline
Question ou thématique associée : à compléter. \\ \hline
\textbf{Étapes de test} &
\begin{enumerate}
    \item Accéder à l'espace administrateur.
    \item Ouvrir la section de gestion des vidéos.
    \item Ajouter ou sélectionner une vidéo.
    \item Choisir la question ou la thématique associée.
    \item Enregistrer l'association.
    \item Vérifier l'affichage côté apprenant.
\end{enumerate}
\\ \hline
\textbf{Résultat attendu} & 
La vidéo est correctement associée au contenu choisi. Elle devient accessible depuis le parcours apprenant attendu. \\ \hline
\textbf{Résultat obtenu} & À compléter après exécution du test. \\ \hline
\textbf{Statut} & Non testé / Validé / Partiellement validé / Non validé \\ \hline
\textbf{Commentaires / anomalies} & À compléter. \\ \hline
\end{tabular}
\caption{Fiche de test A02 -- Association d'une vidéo}
\end{table}

% ============================================================
% Test A02 -- Association d'une vidéo
% ============================================================

%\subsection{Test A02 -- Associer une vidéo à une question ou une thématique}

\begingroup
\small
\renewcommand{\arraystretch}{1.3}
\setlength{\LTleft}{\fill}
\setlength{\LTright}{\fill}

\begin{longtable}{
    |>{\raggedright\arraybackslash}p{0.28\textwidth}
    |>{\raggedright\arraybackslash}p{0.64\textwidth}|
}

\caption{Fiche de test A02 -- Association d'une vidéo, complétée le 13/07/2026}
\label{tab:test-a02-association-video}
\\

\hline
\rowcolor{orange!50}
\multicolumn{2}{|c|}{
    \textbf{Fiche de test A02 -- Association d'une vidéo}
}
\\
\hline

\endfirsthead

\multicolumn{2}{c}{
    \small\textit{Fiche de test A02 -- Association d'une vidéo (suite)}
}
\\
\hline

\rowcolor{orange!50}
\multicolumn{2}{|c|}{
    \textbf{Fiche de test A02 -- Association d'une vidéo}
}
\\
\hline

\endhead

\hline
\multicolumn{2}{|r|}{
    \small\textit{Suite de la fiche page suivante}
}
\\
\hline

\endfoot

\hline
\endlastfoot

% ------------------------------------------------------------
% Informations générales
% ------------------------------------------------------------

\textbf{Identifiant du test}
&
A02
\\
\hline

\textbf{Titre du test}
&
Associer une vidéo pédagogique à une question ou à une thématique.
\\
\hline

\textbf{Type de test}
&
Test fonctionnel de l'interface d'administration et test de liaison entre
les vidéos et les questions.
\\
\hline

\textbf{Persona concerné}
&
Claire, administratrice et formatrice pour Le Carré Pro.
\\
\hline

\textbf{Interface testée}
&
\begin{itemize}[leftmargin=*,nosep]
    \item Interface web d'administration ;
    \item rubrique de gestion des vidéos ;
    \item rubrique de gestion des questions ;
    \item interface apprenant permettant de consulter la vidéo.
\end{itemize}
\\
\hline

\textbf{Test lié}
&
Test A01 -- Création d'une question.

Ce test permet également de contrôler à nouveau l'anomalie
\textbf{A01-04 -- Association des vidéos}.
\\
\hline

\textbf{Date d'exécution}
&
À compléter lors de l'exécution du test.
\\
\hline

\textbf{Exécutant}
&
À compléter.
\\
\hline

\textbf{Environnement}
&
À compléter : environnement de développement, de recette ou de production.
\\
\hline

\textbf{Objectif du test}
&
Vérifier qu'un administrateur peut ajouter ou sélectionner une vidéo
pédagogique, l'associer à une question ou à une thématique existante,
enregistrer cette association et retrouver la vidéo depuis le parcours
apprenant correspondant.
\\
\hline

% ------------------------------------------------------------
% Préconditions
% ------------------------------------------------------------

\rowcolor{orange!50}
\multicolumn{2}{|c|}{
    \textbf{Préconditions}
}
\\
\hline

\textbf{Préconditions générales}
&
\begin{itemize}[leftmargin=*,nosep]
    \item L'application est disponible.
    \item La base de données est accessible.
    \item L'administrateur possède un compte valide.
    \item Le compte utilisé dispose des droits d'administration.
    \item L'administrateur est connecté à l'espace d'administration.
    \item La rubrique de gestion des vidéos est accessible.
    \item La rubrique de gestion des questions est accessible.
    \item Une connexion Internet est disponible pour accéder à YouTube.
\end{itemize}
\\
\hline

\textbf{Contenu nécessaire}
&
\begin{itemize}[leftmargin=*,nosep]
    \item Une question existe déjà dans l'application.
    \item La question créée pendant le test A01 peut être utilisée.
    \item Une thématique permettant de classer la vidéo existe.
    \item Le formulaire permet de saisir une URL vidéo ou de sélectionner
    une vidéo déjà enregistrée.
\end{itemize}
\\
\hline

% ------------------------------------------------------------
% Données de test
% ------------------------------------------------------------

\rowcolor{orange!50}
\multicolumn{2}{|c|}{
    \textbf{Données utilisées pour le test}
}
\\
\hline

\textbf{Référence interne}
&
TEST-A02-VIDEO-ETAT-LIEUX
\\
\hline

\textbf{Titre saisi}
&
État des lieux -- Les pièges à éviter.
\\
\hline

\textbf{Type de ressource}
&
Vidéo externe hébergée sur YouTube.
\\
\hline

\textbf{URL de la vidéo}
&
\url{https://www.youtube.com/watch?v=6xofrG5uZgY&t=24s}
\\
\hline

\textbf{Description proposée}
&
Vidéo pédagogique présentant plusieurs conseils et points de vigilance pour
réaliser un état des lieux et éviter les erreurs susceptibles d'avoir des
conséquences sur la restitution du dépôt de garantie.
\\
\hline

\textbf{Question à associer}
&
Quel est l'objectif principal d'un état des lieux d'entrée ?
\\
\hline

\textbf{Réponse correcte associée}
&
Décrire précisément l'état du local au début de la location.
\\
\hline

\textbf{Type de bail attendu}
&
Bail commercial, lorsque cette valeur est disponible.
\\
\hline

\textbf{Thématique attendue}
&
État des lieux, lorsque cette valeur est disponible.
\\
\hline

\textbf{Valeurs de remplacement}
&
Lorsque les valeurs attendues ne sont toujours pas disponibles, utiliser
les mêmes valeurs que pendant le test A01 afin de poursuivre le test :

\begin{itemize}[leftmargin=*,nosep]
    \item catégorie ou type : \textit{Partie au contrat} ;
    \item thématique : \textit{Professionnelle}.
\end{itemize}

Cette substitution doit être signalée dans le résultat obtenu.
\\
\hline

\textbf{Statut de publication}
&
Vidéo active ou publiée, lorsque ce champ est disponible.
\\
\hline

% ------------------------------------------------------------
% Déroulement
% ------------------------------------------------------------

\rowcolor{orange!50}
\multicolumn{2}{|c|}{
    \textbf{Déroulement du test}
}
\\
\hline

\textbf{Étape 1}
&
Se connecter à l'espace administrateur avec un compte possédant les droits
nécessaires.
\\
\hline

\textbf{Étape 2}
&
Accéder à la rubrique \textit{Vidéos}, \textit{Gestion des vidéos} ou à la
rubrique équivalente.
\\
\hline

\textbf{Étape 3}
&
Noter le nombre de vidéos déjà présentes avant l'ajout de la vidéo de test.
\\
\hline

\textbf{Étape 4}
&
Rechercher la vidéo de test dans la liste afin de vérifier qu'elle n'existe
pas déjà.
\\
\hline

\textbf{Étape 5}
&
Cliquer sur le bouton permettant d'ajouter ou de créer une nouvelle vidéo.
\\
\hline

\textbf{Étape 6}
&
Saisir le titre suivant :

\textit{« État des lieux -- Les pièges à éviter »}.
\\
\hline

\textbf{Étape 7}
&
Saisir l'URL YouTube indiquée dans les données de test.
\\
\hline

\textbf{Étape 8}
&
Saisir la description pédagogique prévue dans les données de test, lorsque
le formulaire contient un champ de description.
\\
\hline

\textbf{Étape 9}
&
Sélectionner le type de bail \textit{Bail commercial} lorsque cette valeur
est proposée.
\\
\hline

\textbf{Étape 10}
&
Sélectionner la thématique \textit{État des lieux} lorsque cette valeur est
proposée.
\\
\hline

\textbf{Étape 11}
&
Activer ou publier la vidéo lorsque cette option est disponible.
\\
\hline

\textbf{Étape 12}
&
Enregistrer la vidéo.
\\
\hline

\textbf{Étape 13}
&
Vérifier qu'un message de confirmation apparaît et qu'aucune erreur
bloquante n'est affichée.
\\
\hline

\textbf{Étape 14}
&
Rechercher la vidéo dans la liste des vidéos à l'aide de son titre ou de sa
référence interne.
\\
\hline

\textbf{Étape 15}
&
Ouvrir la fiche de la vidéo et vérifier que le titre, l'URL, la description,
la thématique et le statut ont été correctement enregistrés.
\\
\hline

\textbf{Étape 16}
&
Cliquer sur le lien ou sur le bouton de prévisualisation de la vidéo.
\\
\hline

\textbf{Étape 17}
&
Vérifier que le lien ouvre la vidéo attendue et qu'il ne conduit pas vers
une page inexistante ou vers une autre vidéo.
\\
\hline

\textbf{Étape 18}
&
Vérifier, lorsque le lecteur intégré le permet, que la lecture commence à
l'instant prévu par l'URL, soit environ 24 secondes après le début.
\\
\hline

% ------------------------------------------------------------
% Association depuis la question
% ------------------------------------------------------------

\rowcolor{orange!50}
\multicolumn{2}{|c|}{
    \textbf{Association de la vidéo à une question}
}
\\
\hline

\textbf{Étape 19}
&
Accéder à la rubrique de gestion des questions.
\\
\hline

\textbf{Étape 20}
&
Rechercher la question suivante :

\textit{« Quel est l'objectif principal d'un état des lieux d'entrée ? »}
\\
\hline

\textbf{Étape 21}
&
Ouvrir la question en mode consultation ou modification.
\\
\hline

\textbf{Étape 22}
&
Localiser le champ permettant d'associer une vidéo à la question.
\\
\hline

\textbf{Étape 23}
&
Ouvrir la liste déroulante des vidéos disponibles.
\\
\hline

\textbf{Étape 24}
&
Vérifier que la vidéo
\textit{« État des lieux -- Les pièges à éviter »}
apparaît dans la liste.

Cette étape permet de vérifier la correction de l'anomalie A01-04.
\\
\hline

\textbf{Étape 25}
&
Sélectionner la vidéo dans la liste.
\\
\hline

\textbf{Étape 26}
&
Enregistrer la modification de la question.
\\
\hline

\textbf{Étape 27}
&
Vérifier qu'un message confirme l'enregistrement de l'association.
\\
\hline

\textbf{Étape 28}
&
Fermer la fiche de la question, puis l'ouvrir à nouveau.
\\
\hline

\textbf{Étape 29}
&
Vérifier que la vidéo précédemment sélectionnée est toujours associée à la
question.
\\
\hline

\textbf{Étape 30}
&
Actualiser complètement la page et vérifier une nouvelle fois que
l'association est conservée.
\\
\hline

\textbf{Étape 31}
&
Se déconnecter, se reconnecter à l'espace administrateur, puis vérifier que
l'association est toujours enregistrée.
\\
\hline

% ------------------------------------------------------------
% Association depuis la vidéo
% ------------------------------------------------------------

\rowcolor{orange!50}
\multicolumn{2}{|c|}{
    \textbf{Vérification depuis la fiche de la vidéo}
}
\\
\hline

\textbf{Étape 32}
&
Retourner dans la rubrique de gestion des vidéos.
\\
\hline

\textbf{Étape 33}
&
Ouvrir la fiche de la vidéo ajoutée pendant le test.
\\
\hline

\textbf{Étape 34}
&
Vérifier si la fiche indique la question ou la thématique à laquelle la
vidéo est associée.
\\
\hline

\textbf{Étape 35}
&
Lorsque l'interface permet d'associer directement une question depuis la
fiche de la vidéo, vérifier que la question sélectionnée est identique à
celle indiquée dans la fiche de la question.
\\
\hline

\textbf{Étape 36}
&
Vérifier qu'il n'existe pas d'incohérence entre l'association affichée dans
la rubrique des vidéos et celle affichée dans la rubrique des questions.
\\
\hline

% ------------------------------------------------------------
% Vérification apprenant
% ------------------------------------------------------------

\rowcolor{orange!50}
\multicolumn{2}{|c|}{
    \textbf{Vérification côté apprenant}
}
\\
\hline

\textbf{Étape 37}
&
Ouvrir l'application avec un compte apprenant autorisé à accéder au contenu
testé.
\\
\hline

\textbf{Étape 38}
&
Accéder au quiz, à la question ou à la thématique utilisée pendant le test.
\\
\hline

\textbf{Étape 39}
&
Répondre à la question afin d'atteindre l'écran de correction ou de
résultat.
\\
\hline

\textbf{Étape 40}
&
Vérifier que la vidéo associée est visible à l'emplacement prévu dans le
parcours apprenant.
\\
\hline

\textbf{Étape 41}
&
Vérifier que le titre affiché est compréhensible et qu'il ne contient pas
de nom technique issu directement de la base de données.
\\
\hline

\textbf{Étape 42}
&
Cliquer sur la vidéo ou sur le lien permettant de la consulter.
\\
\hline

\textbf{Étape 43}
&
Vérifier que la bonne vidéo est ouverte et que la lecture est possible.
\\
\hline

\textbf{Étape 44}
&
Revenir dans l'application et vérifier que le retour au quiz ou au parcours
apprenant fonctionne correctement.
\\
\hline

\textbf{Étape 45}
&
Vérifier que la vidéo n'est pas affichée sur une question ou une thématique
à laquelle elle n'a pas été associée.
\\
\hline

% ------------------------------------------------------------
% Test de modification et de suppression
% ------------------------------------------------------------

\rowcolor{orange!50}
\multicolumn{2}{|c|}{
    \textbf{Modification et suppression de l'association}
}
\\
\hline

\textbf{Étape 46}
&
Retourner dans l'interface administrateur et ouvrir la question utilisée
pendant le test.
\\
\hline

\textbf{Étape 47}
&
Retirer temporairement l'association avec la vidéo, lorsque l'interface le
permet.
\\
\hline

\textbf{Étape 48}
&
Enregistrer la modification.
\\
\hline

\textbf{Étape 49}
&
Vérifier côté apprenant que la vidéo n'est plus affichée pour cette
question.
\\
\hline

\textbf{Étape 50}
&
Associer de nouveau la vidéo et vérifier que celle-ci réapparaît après
l'enregistrement.
\\
\hline

\textbf{Étape 51}
&
À la fin du test, supprimer ou désactiver la vidéo de test et retirer
l'association lorsque ces données ne doivent pas être conservées.
\\
\hline

% ------------------------------------------------------------
% Éléments de preuve
% ------------------------------------------------------------

\rowcolor{orange!50}
\multicolumn{2}{|c|}{
    \textbf{Éléments de preuve à conserver}
}
\\
\hline

\textbf{Captures attendues}
&
\begin{itemize}[leftmargin=*,nosep]
    \item formulaire de création de la vidéo ;
    \item message confirmant son enregistrement ;
    \item vidéo visible dans la liste des vidéos ;
    \item vidéo visible dans le menu déroulant de la question ;
    \item association enregistrée dans la fiche de la question ;
    \item affichage de la vidéo côté apprenant ;
    \item lecture ou ouverture correcte de la vidéo ;
    \item éventuel message d'erreur ou comportement incohérent.
\end{itemize}
\\
\hline

% ------------------------------------------------------------
% Résultat attendu
% ------------------------------------------------------------

\rowcolor{orange!50}
\multicolumn{2}{|c|}{
    \textbf{Résultats du test}
}
\\
\hline

\textbf{Résultat attendu}
&
\begin{itemize}[leftmargin=*,nosep]
    \item L'administrateur peut créer ou enregistrer une vidéo à partir
    d'une URL YouTube valide.

    \item La vidéo apparaît dans la liste de gestion des vidéos.

    \item Les informations saisies sont conservées après actualisation de
    la page.

    \item La vidéo apparaît dans le menu de sélection disponible depuis la
    fiche d'une question.

    \item L'administrateur peut associer la vidéo à la question choisie.

    \item L'association est conservée après fermeture de la fiche,
    actualisation de la page et reconnexion.

    \item La même association est visible depuis les rubriques
    \textit{Questions} et \textit{Vidéos}, lorsque les deux interfaces
    présentent cette information.

    \item La vidéo est affichée uniquement dans le parcours apprenant
    correspondant à la question ou à la thématique choisie.

    \item Le lien ouvre la vidéo attendue.

    \item Le retrait de l'association supprime correctement la vidéo du
    contenu apprenant concerné.

    \item Aucun message contradictoire n'indique qu'aucune vidéo n'est
    disponible lorsque la vidéo existe réellement.
\end{itemize}
\\
\hline

\textbf{Résultat obtenu}
&
À compléter après l'exécution du test.

Indiquer notamment :
\begin{itemize}[leftmargin=*,nosep]
    \item si la vidéo a pu être créée ;
    \item si elle apparaît dans la liste des vidéos ;
    \item si elle apparaît dans la liste déroulante des questions ;
    \item si l'association est conservée après actualisation ;
    \item si la vidéo est accessible côté apprenant ;
    \item si le retrait de l'association fonctionne ;
    \item si l'anomalie A01-04 est toujours présente.
\end{itemize}
\\
\hline

\textbf{Statut}
&
Non testé / Validé / Partiellement validé / Non validé.
\\
\hline

\textbf{Critère de validation}
&
Le test est considéré comme validé lorsque la vidéo peut être enregistrée,
sélectionnée, associée, retrouvée après actualisation et consultée depuis
le parcours apprenant attendu.
\\
\hline

\textbf{Critère de validation partielle}
&
Le test est considéré comme partiellement validé lorsque l'association est
enregistrée mais qu'une incohérence non bloquante subsiste, par exemple :

\begin{itemize}[leftmargin=*,nosep]
    \item intitulé incorrect ;
    \item vidéo absente d'un des écrans d'administration ;
    \item prévisualisation indisponible ;
    \item association visible côté administrateur mais mal présentée côté
    apprenant ;
    \item temps de départ de 24 secondes non pris en compte.
\end{itemize}
\\
\hline

\textbf{Critère de non-validation}
&
Le test est considéré comme non validé lorsque la vidéo ne peut pas être
enregistrée, n'apparaît pas dans la liste de sélection, n'est pas réellement
associée à la question ou reste inaccessible côté apprenant.
\\
\hline

% ------------------------------------------------------------
% Anomalies et observations
% ------------------------------------------------------------

\rowcolor{orange!50}
\multicolumn{2}{|c|}{
    \textbf{Commentaires, observations et anomalies constatées}
}
\\
\hline

\textbf{A02-01 -- Création de la vidéo}
&
La création de la vidéo fonctionne correctement. Les informations saisies
ont été enregistrées dans la base de données sans erreur bloquante.
\\
\hline

\textbf{A02-02 -- Liste des vidéos}
&
Au moment de l'exécution du test, une seule vidéo de test était présente
dans l'application : celle créée dans le cadre de ce scénario.

Elle a donc pu être retrouvée facilement dans la liste. Toutefois, ce test
ne permet pas encore d'évaluer précisément l'efficacité des fonctions de
recherche, de filtrage ou de classement lorsqu'un nombre plus important de
vidéos sera disponible.
\\
\hline

\textbf{A02-03 -- Cohérence de l'association}
&
La vidéo peut être associée à une question et dissociée de celle-ci sans
difficulté technique.

L'association peut être gérée aussi bien depuis la fiche de la question que
depuis la fiche de la vidéo. Les modifications réalisées depuis l'une de
ces interfaces sont correctement répercutées dans l'autre.
\\
\hline

\textbf{A02-04 -- Affichage côté apprenant}
&
L'affichage de la vidéo côté apprenant n'a pas pu être vérifié au cours de
ce test, car l'application apprenant n'était pas accessible dans
l'environnement utilisé.

Cette partie du scénario devra être exécutée ultérieurement depuis
l'application apprenant afin de vérifier :
\begin{itemize}[leftmargin=*,nosep]
    \item la présence de la vidéo dans le parcours attendu ;
    \item l'affichage de son titre ;
    \item l'ouverture correcte du lien ;
    \item la lecture de la vidéo ;
    \item l'absence de la vidéo sur les questions auxquelles elle n'est pas
    associée.
\end{itemize}
\\
\hline

\textbf{A02-05 -- Persistance}
&
L'association est correctement conservée après l'actualisation de la page.

La même vidéo reste associée à la question après le rechargement de
l'interface, ce qui confirme la persistance de l'enregistrement dans la
base de données.
\\
\hline

\textbf{A02-06 -- Association et dissociation}
&
La vidéo peut être associée ou dissociée d'une question depuis deux points
d'entrée :
\begin{itemize}[leftmargin=*,nosep]
    \item depuis la fiche de la question ;
    \item depuis la fiche de la vidéo.
\end{itemize}

La fonctionnalité est opérationnelle dans les deux cas. La gestion depuis
la fiche de la vidéo paraît néanmoins plus simple à comprendre et à
utiliser que celle proposée depuis la fiche de la question.
\\
\hline

\textbf{A02-07 -- Compréhension de l'interface}
&
Le comportement considéré initialement comme une anomalie lors du test A01
correspond en réalité à un fonctionnement peu intuitif de l'interface.

L'association entre une vidéo et une question est bien possible, mais le
parcours utilisateur, les intitulés ou le comportement du menu déroulant ne
permettent pas de comprendre immédiatement la manière dont cette
association doit être réalisée.

Il ne s'agit donc pas d'une anomalie fonctionnelle bloquante, mais d'un
problème d'ergonomie. Une amélioration de l'interface serait souhaitable,
par exemple en :
\begin{itemize}[leftmargin=*,nosep]
    \item ajoutant un libellé plus explicite au champ d'association ;
    \item affichant clairement la vidéo actuellement associée ;
    \item distinguant l'absence d'association de l'absence de vidéos
    disponibles ;
    \item ajoutant un message d'aide ou une indication sur le fonctionnement
    du menu ;
    \item harmonisant le parcours entre la fiche d'une question et celle
    d'une vidéo.
\end{itemize}
\\
\hline

\textbf{Conclusion du test}
&
La création d'une vidéo ainsi que son association et sa dissociation avec
une question sont fonctionnelles. Les données sont correctement
enregistrées et conservées après actualisation de la page.

Le test ne peut toutefois pas être entièrement validé, car l'affichage et
la lecture de la vidéo côté apprenant n'ont pas pu être vérifiés.

Par ailleurs, le mécanisme d'association manque de clarté et peut faire
penser, lors de la première utilisation, que la fonctionnalité ne fonctionne
pas. Une amélioration ergonomique est recommandée afin de rendre le parcours
plus explicite.
\\
\hline

\textbf{Statut proposé}
&
\textbf{Partiellement validé}

La partie administrateur est validée, mais la vérification côté apprenant
reste à effectuer.
\\


\end{longtable}

\endgroup

\subsection{Test A03 -- Consulter les statistiques d'engagement}

\begin{table}[H]
\centering
\renewcommand{\arraystretch}{1.3}
\small
\begin{tabular}{|p{4cm}|p{10cm}|}
\hline
\rowcolor{orange!20}
\multicolumn{2}{|c|}{\textbf{Fiche de test A03 -- Consultation des statistiques}} \\
\hline
\textbf{Identifiant du test} & A03 \\ \hline
\textbf{Titre du test} & Consulter les statistiques d'engagement des apprenants \\ \hline
\textbf{Persona concerné} & Claire, administratrice / formatrice côté Le Carré Pro \\ \hline
\textbf{Interface testée} & Interface administrateur web \\ \hline
\textbf{Objectif du test} & Vérifier que l'administrateur peut consulter des indicateurs utiles au suivi pédagogique et à l'amélioration des contenus. \\ \hline
\textbf{Préconditions} &
\begin{itemize}
    \item L'administrateur est connecté.
    \item Des utilisateurs ont réalisé des quiz.
    \item Des résultats sont disponibles dans la base de données.
\end{itemize}
\\ \hline
\textbf{Données de test} & 
Période d'analyse : à compléter. \newline
Utilisateurs ou groupe test : à compléter. \\ \hline
\textbf{Étapes de test} &
\begin{enumerate}
    \item Accéder à l'espace administrateur.
    \item Ouvrir le tableau de bord ou la section statistiques.
    \item Consulter les indicateurs disponibles.
    \item Vérifier l'affichage du nombre de quiz réalisés.
    \item Vérifier l'affichage du taux de réussite.
    \item Vérifier l'identification des questions ou thématiques difficiles.
\end{enumerate}
\\ \hline
\textbf{Résultat attendu} & 
Les statistiques sont accessibles, lisibles et cohérentes avec les données de test. L'administrateur peut identifier les tendances principales. \\ \hline
\textbf{Résultat obtenu} & À compléter après exécution du test. \\ \hline
\textbf{Statut} & Non testé / Validé / Partiellement validé / Non validé \\ \hline
\textbf{Commentaires / anomalies} & À compléter. \\ \hline
\end{tabular}
\caption{Fiche de test A03 -- Consultation des statistiques}
\end{table}

% ============================================================
% Test A03 -- Consultation des statistiques
% ============================================================

%\subsection{Test A03 -- Consulter les statistiques d'engagement - test du 13/07/2026}

\begingroup
\small
\renewcommand{\arraystretch}{1.3}
\setlength{\LTleft}{\fill}
\setlength{\LTright}{\fill}

\begin{longtable}{
    |>{\raggedright\arraybackslash}p{0.28\textwidth}
    |>{\raggedright\arraybackslash}p{0.64\textwidth}|
}

\caption{Fiche de test A03 -- Consultation des statistiques, complétée le 13/07/2026}
\label{tab:test-a03-statistiques}
\\

\hline
\rowcolor{orange!50}
\multicolumn{2}{|c|}{
    \textbf{Fiche de test A03 -- Consultation des statistiques}
}
\\
\hline

\endfirsthead

\multicolumn{2}{c}{
    \small\textit{Fiche de test A03 -- Consultation des statistiques (suite)}
}
\\
\hline

\rowcolor{orange!50}
\multicolumn{2}{|c|}{
    \textbf{Fiche de test A03 -- Consultation des statistiques}
}
\\
\hline

\endhead

\hline
\multicolumn{2}{|r|}{
    \small\textit{Suite de la fiche page suivante}
}
\\
\hline

\endfoot

\hline
\endlastfoot

% ------------------------------------------------------------
% Informations générales
% ------------------------------------------------------------

\textbf{Identifiant du test}
&
A03
\\
\hline

\textbf{Titre du test}
&
Consulter les statistiques d'engagement des apprenants.
\\
\hline

\textbf{Type de test}
&
Test fonctionnel, test de cohérence des données et test d'ergonomie.
\\
\hline

\textbf{Persona concerné}
&
Claire, administratrice et formatrice pour Le Carré Pro.
\\
\hline

\textbf{Interface testée}
&
Interface web d'administration, rubrique
\textit{Monitoring des statistiques}.
\\
\hline

\textbf{Date d'exécution}
&
À compléter.
\\
\hline

\textbf{Exécutant}
&
À compléter.
\\
\hline

\textbf{Objectif du test}
&
Vérifier que l'administrateur peut consulter des indicateurs utiles au suivi
pédagogique, contrôler leur cohérence et identifier les questions présentant
un faible taux de réussite.

Le test doit également vérifier le comportement de l'interface lorsque
certaines questions ou vidéos ne disposent encore d'aucune donnée
d'utilisation.
\\
\hline

% ------------------------------------------------------------
% Préconditions
% ------------------------------------------------------------

\rowcolor{orange!50}
\multicolumn{2}{|c|}{
    \textbf{Préconditions}
}
\\
\hline

\textbf{Préconditions générales}
&
\begin{itemize}[leftmargin=*,nosep]
    \item L'application et la base de données sont accessibles.
    \item L'administrateur possède un compte valide.
    \item Le compte utilisé dispose des droits d'administration.
    \item L'administrateur est connecté à l'espace d'administration.
    \item La rubrique de statistiques est accessible.
    \item Plusieurs réponses à des questions sont déjà enregistrées.
\end{itemize}
\\
\hline

\textbf{Limites du jeu de données}
&
Le volume de données disponible au moment du test reste limité.

Les résultats permettent de vérifier le fonctionnement et la cohérence
générale de l'interface, mais ils ne sont pas suffisants pour réaliser une
analyse pédagogique représentative de l'ensemble des futurs apprenants.
\\
\hline

% ------------------------------------------------------------
% Données de référence
% ------------------------------------------------------------

\rowcolor{orange!50}
\multicolumn{2}{|c|}{
    \textbf{Données observées avant le test}
}
\\
\hline

\textbf{Nombre total affiché}
&
197 questions jouées ou occurrences de réponses enregistrées.
\\
\hline

\textbf{Bonnes réponses}
&
41 bonnes réponses.
\\
\hline

\textbf{Taux de réussite global}
&
21\,\%.
\\
\hline

\textbf{Temps moyen global}
&
1,5 seconde par réponse.
\\
\hline

\textbf{Vues de vidéos}
&
0 vue.
\\
\hline

\textbf{Taux de complétion vidéo}
&
0\,\%.
\\
\hline

\textbf{Questions présentes dans la liste statistique}
&
146 questions disposant apparemment d'au moins une donnée d'utilisation.
\\
\hline

\textbf{Question de contrôle}
&
Question \#466 :

\textit{« Parmi ces exemples, lequel justifie typiquement une convention
d'occupation précaire ? »}
\\
\hline

\textbf{Données de contrôle de la question \#466}
&
\begin{itemize}[leftmargin=*,nosep]
    \item catégorie : \textit{Fondamentaux} ;
    \item type de bail : \textit{Précaire} ;
    \item niveau : niveau 1 ;
    \item parties jouées : 5 ;
    \item bonnes réponses : 1 ;
    \item mauvaises réponses : 4 ;
    \item taux de réussite : 20\,\% ;
    \item temps cumulé : 5,8 secondes ;
    \item temps moyen : 1,2 seconde ;
    \item dernière partie : 8 juin 2026 à 15 h 00.
\end{itemize}
\\
\hline

\textbf{Question sans statistique}
&
La question créée lors du test A01 ne figure pas dans la liste, car elle
n'a encore jamais été jouée.
\\
\hline

% ------------------------------------------------------------
% Déroulement
% ------------------------------------------------------------

\rowcolor{orange!50}
\multicolumn{2}{|c|}{
    \textbf{Déroulement du test}
}
\\
\hline

\textbf{Étape 1}
&
Se connecter à l'espace administrateur avec un compte disposant des droits
nécessaires.
\\
\hline

\textbf{Étape 2}
&
Accéder à la rubrique \textit{Monitoring des statistiques}.
\\
\hline

\textbf{Étape 3}
&
Vérifier que la page se charge sans erreur et que les indicateurs globaux
sont visibles.
\\
\hline

\textbf{Étape 4}
&
Relever les valeurs affichées pour :
\begin{itemize}[leftmargin=*,nosep]
    \item le nombre de questions jouées ;
    \item le nombre de bonnes réponses ;
    \item le taux de réussite global ;
    \item le temps moyen par réponse ;
    \item le nombre de vues de vidéos ;
    \item le taux de complétion des vidéos.
\end{itemize}
\\
\hline

\textbf{Étape 5}
&
Vérifier la cohérence du taux de réussite global à partir des données
affichées.

Le calcul attendu est :

\[
\frac{41}{197} \times 100 \approx 20{,}8\,\%
\]

La valeur affichée peut donc être arrondie à 21\,\%.
\\
\hline

\textbf{Étape 6}
&
Vérifier que la liste des statistiques par question est affichée.
\\
\hline

\textbf{Étape 7}
&
Vérifier la présence des colonnes suivantes :
\begin{itemize}[leftmargin=*,nosep]
    \item identifiant ;
    \item question ;
    \item type de bail ;
    \item nombre de parties jouées ;
    \item taux de réussite ;
    \item temps moyen de réponse ;
    \item accès au détail.
\end{itemize}
\\
\hline

\textbf{Étape 8}
&
Utiliser le champ de recherche pour retrouver la question \#466 à partir
de son identifiant ou d'une partie de son énoncé.
\\
\hline

\textbf{Étape 9}
&
Vérifier que le résultat correspondant est correctement affiché.
\\
\hline

\textbf{Étape 10}
&
Tester successivement les filtres disponibles :
\begin{itemize}[leftmargin=*,nosep]
    \item catégorie ;
    \item type de bail ;
    \item statut.
\end{itemize}
\\
\hline

\textbf{Étape 11}
&
Combiner plusieurs filtres et vérifier que la liste est correctement mise
à jour.
\\
\hline

\textbf{Étape 12}
&
Réinitialiser les filtres et vérifier que la liste complète réapparaît.
\\
\hline

\textbf{Étape 13}
&
Effectuer une recherche ne correspondant à aucune question et vérifier
qu'un message explicite indique l'absence de résultat.
\\
\hline

\textbf{Étape 14}
&
Cliquer sur le bouton \textit{Inspecter} de la question \#466.
\\
\hline

\textbf{Étape 15}
&
Vérifier que la fiche de détail affiche correctement :
\begin{itemize}[leftmargin=*,nosep]
    \item l'énoncé ;
    \item la catégorie ;
    \item le type de bail ;
    \item le niveau ;
    \item la date de dernière partie ;
    \item les parties jouées ;
    \item les bonnes et mauvaises réponses ;
    \item le taux de réussite ;
    \item le temps cumulé ;
    \item le temps moyen par réponse.
\end{itemize}
\\
\hline

\textbf{Étape 16}
&
Vérifier la cohérence du taux de réussite de la question \#466 :

\[
\frac{1}{5} \times 100 = 20\,\%
\]
\\
\hline

\textbf{Étape 17}
&
Vérifier la cohérence du temps moyen de réponse :

\[
\frac{5{,}8}{5} = 1{,}16\ \text{seconde}
\]

La valeur peut être arrondie à 1,2 seconde.
\\
\hline

\textbf{Étape 18}
&
Fermer la fiche de détail et vérifier que l'utilisateur revient à la liste
sans perdre sa recherche ou ses filtres, lorsque ce comportement est prévu.
\\
\hline

\textbf{Étape 19}
&
Rechercher la question créée pendant le test A01.
\\
\hline

\textbf{Étape 20}
&
Constater si la question apparaît dans la liste alors qu'elle n'a jamais
été jouée.
\\
\hline

\textbf{Étape 21}
&
Lorsque l'application apprenant est disponible, jouer une fois la question
créée pendant le test A01.
\\
\hline

\textbf{Étape 22}
&
Actualiser la page des statistiques et vérifier que la question apparaît
après sa première utilisation.
\\
\hline

\textbf{Étape 23}
&
Ouvrir l'onglet \textit{Vidéos}.
\\
\hline

\textbf{Étape 24}
&
Vérifier que l'interface affiche un état vide compréhensible lorsqu'aucune
vidéo n'a encore été visionnée.
\\
\hline

\textbf{Étape 25}
&
Actualiser complètement la page et vérifier que les statistiques restent
accessibles et identiques.
\\
\hline

\textbf{Étape 26}
&
Conserver des captures d'écran du tableau de bord, de la liste filtrée et
de la fiche détaillée de la question \#466.
\\
\hline

% ------------------------------------------------------------
% Résultat attendu
% ------------------------------------------------------------

\rowcolor{orange!50}
\multicolumn{2}{|c|}{
    \textbf{Résultats du test}
}
\\
\hline

\textbf{Résultat attendu}
&
\begin{itemize}[leftmargin=*,nosep]
    \item La rubrique de statistiques est accessible sans erreur.
    \item Les indicateurs globaux sont lisibles.
    \item Les calculs affichés sont cohérents avec les données enregistrées.
    \item La liste des questions peut être recherchée et filtrée.
    \item Le détail d'une question est accessible.
    \item Les statistiques détaillées sont cohérentes avec les valeurs de
    la liste.
    \item Les questions sans activité restent identifiables avec des valeurs
    nulles, ou un filtre permet de les afficher.
    \item L'absence de données vidéo est présentée avec un message clair.
    \item Les données sont conservées après actualisation de la page.
\end{itemize}
\\
\hline

\textbf{Résultat obtenu}
&
\begin{itemize}[leftmargin=*,nosep]
    \item La page de monitoring est accessible et les principaux indicateurs
    sont affichés.

    \item Le tableau de bord indique 197 occurrences de questions jouées,
    41 bonnes réponses, un taux de réussite global de 21\,\% et un temps
    moyen de 1,5 seconde par réponse.

    \item Le calcul du taux de réussite global est cohérent avec les données
    affichées.

    \item La liste statistique contient 146 questions disposant de données
    d'utilisation.

    \item La fiche de détail de la question \#466 est accessible.

    \item Les valeurs détaillées de cette question sont cohérentes :
    5 parties jouées, 1 bonne réponse, 4 mauvaises réponses et un taux de
    réussite de 20\,\%.

    \item Le temps cumulé de 5,8 secondes correspond à un temps moyen arrondi
    à 1,2 seconde.

    \item Aucune vue de vidéo et aucune lecture complète ne sont actuellement
    enregistrées.

    \item La question créée lors du test A01 n'apparaît pas dans les
    statistiques, car elle n'a jamais été jouée.

    \item Les tests de filtrage, de recherche sans résultat et d'affichage
    côté vidéo restent à compléter selon les manipulations réalisées.
\end{itemize}
\\
\hline

\textbf{Statut proposé}
&
\textbf{Partiellement validé}
\\
\hline

\textbf{Justification du statut}
&
Les principales statistiques sont accessibles et les calculs vérifiés sont
cohérents.

Le test reste partiellement validé en raison du faible volume de données,
de l'absence de statistiques vidéo et de l'impossibilité d'identifier
directement les questions qui n'ont jamais été jouées.
\\
\hline

% ------------------------------------------------------------
% Observations et anomalies
% ------------------------------------------------------------

\rowcolor{orange!50}
\multicolumn{2}{|c|}{
    \textbf{Commentaires, observations et anomalies constatées}
}
\\
\hline

\textbf{A03-01 -- Cohérence des indicateurs globaux}
&
Le taux de réussite global affiché est cohérent avec le nombre de bonnes
réponses et le nombre total d'occurrences enregistrées.

Les 41 bonnes réponses sur 197 occurrences représentent environ 20,8\,\%,
soit 21\,\% après arrondi.
\\
\hline

\textbf{A03-02 -- Cohérence du détail d'une question}
&
Les données de la question \#466 sont cohérentes.

Une bonne réponse sur cinq parties jouées correspond bien à un taux de
réussite de 20\,\%. Le temps cumulé de 5,8 secondes correspond également à
un temps moyen arrondi à 1,2 seconde.
\\
\hline

\textbf{A03-03 -- Ambiguïté du libellé « Questions jouées »}
&
Le tableau de bord affiche 197 \textit{questions jouées}, tandis que la
liste contient 146 questions distinctes.

Le compteur de 197 semble donc correspondre au nombre total d'occurrences
ou de réponses enregistrées, et non au nombre de questions distinctes.

Un libellé tel que \textit{« Réponses enregistrées »},
\textit{« Tentatives »} ou \textit{« Occurrences jouées »} serait plus
précis.
\\
\hline

\textbf{A03-04 -- Questions jamais jouées invisibles}
&
Les questions qui n'ont encore jamais été jouées ne figurent pas dans la
liste statistique.

La question créée lors du test A01 est donc absente, bien qu'elle existe
dans la base de données et qu'elle soit disponible dans l'interface de
gestion des questions.

Ce comportement empêche l'administrateur d'identifier facilement les
contenus qui ne sont jamais proposés ou jamais utilisés.

Il serait préférable :
\begin{itemize}[leftmargin=*,nosep]
    \item d'afficher toutes les questions avec la valeur
    \textit{0 partie jouée} ;
    \item ou d'ajouter un filtre
    \textit{« Inclure les questions jamais jouées »} ;
    \item ou de créer un indicateur recensant les contenus sans activité.
\end{itemize}
\\
\hline

\textbf{A03-05 -- Volume de données limité}
&
Le nombre actuel de parties ne permet pas encore de tirer des conclusions
pédagogiques fiables.

Par exemple, une question réussie une fois sur une seule tentative
obtiendrait un taux de réussite de 100\,\%, sans que ce résultat soit
représentatif.

Le nombre de tentatives doit donc toujours être affiché à côté du taux de
réussite.
\\
\hline

\textbf{A03-06 -- Absence de période d'analyse}
&
Aucun filtre de période n'est visible sur l'écran testé.

L'administrateur ne peut donc pas distinguer les résultats de la semaine,
du mois, d'une promotion ou de l'ensemble de la durée d'utilisation.

L'ajout d'un sélecteur de dates permettrait de suivre l'évolution de
l'engagement dans le temps.
\\
\hline

\textbf{A03-07 -- Absence de segmentation des apprenants}
&
Les statistiques affichées sont globales.

Aucun filtre visible ne permet de sélectionner un utilisateur, une
promotion, un groupe ou un type d'abonnement. Cette fonctionnalité serait
utile pour réaliser un suivi pédagogique plus précis.
\\
\hline

\textbf{A03-08 -- Statistiques vidéo non exploitables}
&
Aucune vue de vidéo n'est actuellement enregistrée. Les indicateurs de vues
et de complétion restent donc à zéro.

L'interface devrait afficher un message explicite indiquant qu'aucune donnée
n'est encore disponible, plutôt que de laisser uniquement apparaître des
valeurs nulles.
\\
\hline

\textbf{Conclusion du test}
&
La consultation des statistiques fonctionne et les calculs contrôlés sont
cohérents avec les données affichées.

L'interface permet déjà d'identifier certaines questions présentant un
faible taux de réussite. Elle reste cependant limitée pour une véritable
analyse pédagogique, notamment en raison de l'absence des questions jamais
jouées, du manque de filtres temporels et de segmentation, ainsi que du
faible volume actuel de données.

La fonctionnalité est donc opérationnelle, mais plusieurs améliorations
permettraient de rendre les statistiques plus complètes et plus utiles pour
l'administrateur.
\\
\hline

\end{longtable}

\endgroup